% Every number below is a macro from tables/numbers.tex (scripts/make_paper_tables.py). Rule outcomes are quoted
% from the same file (\Rule... macros), so a changed outcome shows up in the text rather than contradicting it.

\subsection{Models and detectors}
The honest model explains the test prices well ($R^2$ \HonestRtwoFeature\ in the feature design, \HonestRtwoLocation\
in the location design) and the innocuous model, which sees neither location nor demographics, less well
(\PsiRtwoFeature\ and \PsiRtwoLocation; Table~\ref{tab:data}). The biased model penalises \SharePenFeature\% of test
homes, about \NPenFeature\ of the 200 explained homes in each seed. In the location design the model's coordinate
lookup puts \RasterAgree\% of test homes in the same tract as the spatial join in the data file.

The matched detectors separate LIME's queries from real homes completely (AUC \FLIMEMatAuc), but separate the
Shapley explainers' queries only partly (AUC \FKernelSHAPMatAuc\ for KernelSHAP and \FGeoShapleyMatAuc\ for
GeoShapley in the feature design). This difference, not the attack recipe, drives everything that follows.

\subsection{A scaffold hides the penalty from all three explainers}
\figinput{F2_main}{0.66}{Penalty share kept in the feature design (1 = the explanation shows the whole penalty, as it
does when the biased model is served directly; 0 = none of it). Mean and 95\% interval over seeds.}

\tabinput{tab_main_feature}

\paragraph{Feature design.} Served directly, the biased model's penalty is plain in every explanation: the trigger
is among the three most important features for \FLIMEBiasedTop\% (LIME), \FKernelSHAPBiasedTop\% (KernelSHAP) and
\FGeoShapleyBiasedTop\% (GeoShapley) of the explained homes, with a signed attribution on penalised homes of about
\FKernelSHAPBiasedSigned\ in log$_{10}$ price (Table~\ref{tab:main_feature}).

The generic detector of \citet{slack2020fooling} does not fool LIME here: it lets \FLIMEGenLeak\% of LIME's queries
through to $f$, and LIME keeps \FLIMEGenKept\ of the penalty. This is the source of the earlier project finding that
LIME resists the attack in a geographic setting (Appendix~D); it is a failure of the detector, not a property of
LIME. With a detector matched to LIME's own sampler, no LIME query reaches $f$ (\FLIMEMatLeak\%), the scaffold still
serves $f$ to \FLIMEMatFid\% of the explained homes, and LIME keeps \FLIMEMatKept\ of the penalty, with the trigger in
the top three for \FLIMEMatTop\% of homes. LIME is \RuleConePrimeFeatureLIME.

Against the Shapley explainers the matched scaffold, at \FKernelSHAPMatFid\% and \FGeoShapleyMatFid\% fidelity, cuts
the penalty to \FKernelSHAPMatKept\ (95\% interval \FKernelSHAPMatKeptLo--\FKernelSHAPMatKeptHi) for KernelSHAP and
\FGeoShapleyMatKept\ (\FGeoShapleyMatKeptLo--\FGeoShapleyMatKeptHi) for GeoShapley, and the trigger drops out of the
top three for about half of the homes (\FKernelSHAPMatTop\% and \FGeoShapleyMatTop\% remain). Rule C1 (Holm-corrected)
gives LIME \RuleConeFeatureLIME, KernelSHAP \RuleConeFeatureKernelSHAP\ and GeoShapley \RuleConeFeatureGeoShapley;
by rule C1$'$ KernelSHAP is \RuleConePrimeFeatureKernelSHAP\ and GeoShapley \RuleConePrimeFeatureGeoShapley. The detector cannot do better because \FKernelSHAPMatLeak\% of KernelSHAP's and
\FGeoShapleyMatLeak\% of GeoShapley's queries look real enough to be served by $f$.

\paragraph{Location design.} Here location has a legitimate price effect, so even the honest model's explanation
puts a negative location attribution on the homes the biased model penalises (\LKernelSHAPHonestSigned\ for
KernelSHAP). Table~\ref{tab:main_location} therefore gives the share kept both raw and net of the honest model. The
LIME-matched scaffold removes location from LIME's explanation altogether: the raw share is \LLIMEMatKeptRaw, location
is in the top three for \LLIMEMatTop\% of homes against \LLIMEHonestTop\% under the honest model, and the net share
is negative (\LLIMEMatKept): LIME now shows less location effect than an honest model would. KernelSHAP and
GeoShapley keep \LKernelSHAPMatKept\ and \LGeoShapleyMatKept\ of the penalty net of the honest model (raw
\LKernelSHAPMatKeptRaw\ and \LGeoShapleyMatKeptRaw), at fidelity \LKernelSHAPMatFid\% and \LGeoShapleyMatFid\%.
Location stays in their top three for most homes (\LKernelSHAPMatTop\% and \LGeoShapleyMatTop\%), as it should for a
housing model; what the scaffold removes is about half of the penalty carried by it. In this design rule C1 gives
LIME \RuleConeLocationLIME, KernelSHAP \RuleConeLocationKernelSHAP\ and GeoShapley \RuleConeLocationGeoShapley, and
rule C1$'$ gives LIME \RuleConePrimeLocationLIME, KernelSHAP \RuleConePrimeLocationKernelSHAP\ and GeoShapley
\RuleConePrimeLocationGeoShapley.

\tabinput{tab_main_location}

\subsection{What hiding costs: the stealth--fidelity frontier}
\figinput{F3_frontier}{0.92}{The stealth--fidelity frontier. Each filled marker is the matched detector at a stricter
threshold (90\%, 75\% and 50\% fidelity on the calibration homes); hollow markers are the generic detector at 0.5.
The location panel shows the share kept net of the honest model. Mean and 95\% interval over seeds in both
directions.}

An attacker who wants to hide more from a Shapley explainer has one lever: withhold the bias from more real homes.
Lowering the matched threshold so that 75\% and then 50\% of real homes receive $f$ lowers KernelSHAP's share kept
from \FKernelSHAPMatKept\ to \FKernelSHAPMatSeventyFiveKept\ and \FKernelSHAPMatFiftyKept, and GeoShapley's from
\FGeoShapleyMatKept\ to \FGeoShapleyMatSeventyFiveKept\ and \FGeoShapleyMatFiftyKept\ (Table~\ref{tab:frontier},
Figure~\ref{fig:F3_frontier}). The share kept falls roughly in step with fidelity. Even at 50\% fidelity --- a
scaffold that spares half the homes it was built to penalise --- the trigger is still in the top three for
\FKernelSHAPMatFiftyTop\% (KernelSHAP) and \FGeoShapleyMatFiftyTop\% (GeoShapley) of homes.

The generic detector, trained on different perturbations and run at a fixed threshold, reaches
\FKernelSHAPGenFid\% fidelity and a share kept of \FKernelSHAPGenKept\ (KernelSHAP) and \FGeoShapleyGenKept\
(GeoShapley), on the same curve as the matched detector within the seed intervals. Within the detectors we tried, the
trade-off is set by where the explainer queries the model, not by how the attacker builds its detector. In the
location design the net share falls to \LKernelSHAPMatSeventyFiveKept\ and \LGeoShapleyMatSeventyFiveKept\ at 75\%
fidelity and to about zero at 50\% (\LKernelSHAPMatFiftyKept, \LGeoShapleyMatFiftyKept).

\tabinput{tab_frontier}

\subsection{GeoShapley is fooled like KernelSHAP, and its decomposition changes}
\paragraph{No detectable difference from KernelSHAP (rule C2).} Paired by seed, GeoShapley minus KernelSHAP in the
share kept under the matched scaffold is \FGeoMinusKsKept\ (interval \FGeoMinusKsKeptLo\ to \FGeoMinusKsKeptHi) in the
feature design and \LGeoMinusKsKept\ (\LGeoMinusKsKeptLo\ to \LGeoMinusKsKeptHi) in the location design: feature
design, \RuleCtwoFeatureGeoMinusKs; location design, \RuleCtwoLocationGeoMinusKs. GeoShapley's queries are
somewhat harder to tell from real homes (difference in detector AUC \FGeoMinusKsAuc\ in the feature design), consistent
with its moving both coordinates together so that every hybrid sits at a real location, but this does not translate
into a detectable difference in what the attacker can hide. Treating location as a joint player does not make an
explainer harder to fool; the coalition construction it shares with KernelSHAP decides that.

\paragraph{The decomposition moves (rule C3).} Under the matched scaffold the intrinsic location effect on penalised
homes moves towards zero and the feature-by-location interactions grow (Table~\ref{tab:geo},
Figure~\ref{fig:F4_geoshapley_decomposition}). In the feature design GEO on penalised homes goes from
\FGeoBiasedGeoPen\ under the biased model to \FGeoMatGeoPen\ under the matched scaffold, and the per-home sum of
absolute interaction effects from \FGeoBiasedInt\ to \FGeoMatInt. In the location design GEO goes from
\LGeoBiasedGeoPen\ (honest model: \LGeoHonestGeoPen) to \LGeoMatGeoPen, and the interactions from \LGeoBiasedInt\ to
\LGeoMatInt. Each change has a paired interval that excludes zero. The scaffold's routing depends on how real a hybrid
looks rather than on any single feature, and GeoShapley reports the resulting discontinuities as features acting
differently at different locations. An auditor who reads only the GEO map sees less of the penalty than the location
player as a whole carries.

\figinput{F4_geoshapley_decomposition}{0.82}{GeoShapley's decomposition by deployment: the intrinsic location effect
(GEO) on penalised homes, and the per-home sum of absolute feature-by-location interaction effects. Mean and 95\%
interval over seeds.}

\tabinput{tab_geo}

\subsection{Why: efficiency and coalition sizes}
\paragraph{Efficiency.} Under the matched scaffold the Shapley explanations of penalised homes still carry
\FKernelSHAPMatGapTotal\ (KernelSHAP) and \FGeoShapleyMatGapTotal\ (GeoShapley) of the penalty gap the biased model's
explanation puts between penalised and unpenalised homes: \FKernelSHAPMatKept\ and \FGeoShapleyMatKept\ on the trigger
and \FKernelSHAPMatReloc\ and \FGeoShapleyMatReloc\ moved to other features, most of it to \FKernelSHAPMatRelocTo.
Because the background term in the efficiency identity is the same for every home, the rest can leave the
explanation only through explained homes that the scaffold itself serves with $\psi$. LIME's explanation carries
\FLIMEMatGapTotal\ of the gap. The penalty that efficiency keeps in a Shapley explanation can be moved but
not removed while the explained home is served by $f$.

\paragraph{Routing by coalition size.} Figure~\ref{fig:F7_mechanism} shows the share of KernelSHAP's hybrids routed
to $f$ at each coalition size. It is U-shaped: \MechFeatureKernelSHAPRhoFirst\ for a background home with only the
trigger swapped in, \MechFeatureKernelSHAPRhoLast\ for the explained home itself ($k = \MechFeatureKernelSHAPM$), and
\MechFeatureKernelSHAPRhoMin\ at its lowest ($k = \MechFeatureKernelSHAPRhoMinK$). Both ends of the coalition range
are near-real homes the detector cannot reject without rejecting real homes; only the middle is far enough from the
data to be routed away, and equation~\eqref{eq:sizes} gives every size the same weight. Stricter thresholds lower the
whole curve, the ends included, which is why hiding costs fidelity.

\paragraph{Prediction (rule M).} The independent Monte Carlo evaluation of equation~\eqref{eq:sizes} predicts the
measured share kept with a mean absolute error of \MechFeatureKernelSHAPMAE\ for KernelSHAP (seed-to-seed standard
deviation of the measured share \MechFeatureKernelSHAPSD; \RuleMechFeatureKernelSHAP), \MechFeatureGeoShapleyMAE\ for
GeoShapley in the feature design (\RuleMechFeatureGeoShapley) and \MechLocationGeoShapleyMAE\ in the location design
(\RuleMechLocationGeoShapley; Table~\ref{tab:mechanism}). The plain average of the routing rates over sizes, which
ignores how the penalty is spread across sizes, is within \MechFeatureKernelSHAPRhoMAE\ of the measured share for
KernelSHAP: to a first approximation, \emph{the share of the penalty a Shapley explainer keeps is the average share of
its hybrids the detector cannot reject}.

\figinput{F7_mechanism}{0.92}{Mechanism. Left: share of with-trigger hybrids routed to $f$ by coalition size, at the
three attacker thresholds (feature design). Right: predicted against measured penalty share kept, one point per seed
and threshold; crosses are the plain average of the routing rates.}

\tabinput{tab_mechanism}

\subsection{Robustness}
\paragraph{Model families and audit settings.} Table~\ref{tab:robust} and Figure~\ref{fig:F8_robustness} repeat the
feature-design attack with random forests (honest $R^2$ \HonestRtwoForest) and neural networks (\HonestRtwoNet) in
place of XGBoost, and against auditors who change the explainer's settings after the attacker has trained its
detector. In \RobHold\ of \RobCells\ explainer--configuration cells the C1$'$ class is the same as in the main result
(rule R); exceptions: \RobBreakList. \todo{one or two sentences on the direction of the settings effects once all
seeds are in: does a larger background or $k$-means background raise the share kept, does LIME's Gaussian sampler
escape the quartile-trained detector?}

\tabinput{tab_robust}
\figinput{F8_robustness}{0.85}{Penalty share kept under the matched scaffold across model families and audit
settings the attacker did not train for (feature design). Mean and 95\% interval over seeds.}

\paragraph{Spatial autocorrelation of the trigger.} \todo{C5 result from Table~\ref{tab:autocorr} and
Figure~\ref{fig:F6_autocorrelation}: slope of the matched leak on Moran's $I$ (\RuleCfiveFeature).}
\tabinput{tab_autocorr}
\figinput{F6_autocorrelation}{0.66}{Queries reaching $f$ under the matched scaffold as the trigger's spatial
autocorrelation is lowered. Mean and 95\% interval over seeds.}

\paragraph{A smaller penalty.} \todo{Table~\ref{tab:penalty010}: share kept at $\delta = 0.10$ (LIME \PenTenLIMEKept,
KernelSHAP \PenTenKernelSHAPKept, GeoShapley \PenTenGeoShapleyKept).}
\tabinput{tab_penalty010}

\subsection{An outcome-based check sees what the explainers miss}
No explainer is needed to find this penalty. On the 600 held-out homes, the 8-nearest-neighbour mean of the
deployed model's residuals separates penalised from unpenalised homes with AUC \EsevenScaffoldAucLagThirty\ at
$\delta = 0.30$ and \EsevenAucLagScaffoldTen\ (interval \EsevenAucLagScaffoldTenLo--\EsevenAucLagScaffoldTenHi) at
$\delta = 0.10$, for the scaffold that fools LIME completely (Table~\ref{tab:e7}, Figure~\ref{fig:F5_outcome_check}).
The check loses power only for penalties of a few percent (AUC \EsevenScaffoldAucLagTwo\ at $\delta = 0.02$). It costs
the auditor few false alarms: LISA flags \EsevenHonestFlag\% of the honest model's unpenalised homes as low--low
clusters. The scaffold cannot evade this check without changing what real homes are charged, because it is computed
from served predictions and observed prices only. By rule E7 (\RuleEsevenFeature) the check is reported here; it is a
complement to explanation-based audits, not a defence we evaluate against adaptive attackers.

\figinput{F5_outcome_check}{0.6}{Outcome-based check: AUC of the local mean residual for the penalised homes, by
penalty size, for the biased model and the LIME-matched scaffold; the line marks the honest model. Mean and 95\%
interval over seeds.}

\tabinput{tab_e7}
